\documentclass[10pt,a4paper]{amsart}
\usepackage[margin=19mm]{geometry}
\usepackage{amsmath,amssymb,mathtools}
\usepackage[scr=rsfs,cal=euler]{mathalfa}
\usepackage{xfrac}
\usepackage[unicode,hidelinks,bookmarks=false]{hyperref}
\newcommand{\ie}{i.e.\ }
\newcommand{\cf}{cf.\ }
\newcommand{\resp}{respectively\ }
\newcommand{\Subsection}{\subsection}
\providecommand{\nobreakdash}{\nobreak\hbox{-}}
\setlength{\emergencystretch}{2em}
\multlinegap0pt
\usepackage{bm}
\usepackage{bbm}
\usepackage{dsfont}
\usepackage{xspace}
\usepackage{cancel}
\usepackage{mathtools}
\usepackage{amsthm}
\makeatletter
\@ifundefined{exam}{\newtheorem{exam}{Exam}}{}
\@ifundefined{lem}{\newtheorem{lem}[exam]{Lemma}}{}
\makeatother
\newcommand{\cms}{\operatorname{comass}}
\newcommand{\du}{^\ast}
\newcommand{\ms}{\mathbf{M}}
\newcommand{\no}[1]{\left\lVert#1\right\rVert}
\newcommand{\pf}{_\ast}
\newcommand{\s}{\subset}
\title[On a conjecture of Almgren: area-minimizing submanifolds wit]{On a conjecture of Almgren: area-minimizing submanifolds with fractal singularities}
\author{Zhenhua Liu}
\date{Source: arXiv:2110.13137v7 (CC BY 4.0)}
\begin{document}
\maketitle

\noindent\emph{Source and licence.} On a conjecture of Almgren: area-minimizing submanifolds with fractal singularities. Version arXiv:2110.13137v7 (CC BY 4.0), distributed under the Creative Commons Attribution 4.0 International licence, as recorded in the arXiv licence metadata for this exact version and shown on the abstract page https://arxiv.org/abs/2110.13137v7. Full rights notice: licenses/arxiv-2110.13137-CC-BY-4.0.txt.\par\smallskip

\noindent\textit{Editorial scope.} Counted unit: the Lem whose source label is \texttt{lemprod} in the article (source0.tex lines 715--728, source label \texttt{lemprod}), delivered together with the article's own complete original proof of it (source0.tex lines 729--751, one \texttt{proof} environment of 23 lines). No mathematics is written, repaired or re-derived: statement and proof are the article's own text line for line. Required context retained uncounted: eqcms (lines 401--403). Every definition, display and label of those lines is retained unchanged. Omitted from this package and not redistributed: 1--400, 404--714, 752--1273. Nothing inside a retained excerpt is omitted or altered: every retained body is delivered exactly as the article's lines are written, so no omission receipt of any kind accompanies this package. Citations: the retained excerpts carry no citation command at all, so no bibliography entry of the article is transcribed and no reference is redistributed. Reference adaptation: none was needed and none was made; every reference inside the delivered unit resolves inside it, so no unresolved reference or citation is produced. Printed slips kept literal: no printed word, symbol, hypothesis or number of the counted statement or of its proof was altered. Layout and class: the article's own class is \texttt{amsart} with packages including amsmath, caption, booktabs, lipsum, amsthm, pdfpages, relsize, amsfonts, amssymb, mathrsfs, tikz, environ, elocalloc, url; the wrapper is the lane's pinned \texttt{amsart} editorial wrapper at 10pt on A4, which restates the theorem-like environments the retained text needs and adds the article's own macro definitions that the retained text needs (\texttt{\textbackslash cms}, \texttt{\textbackslash du}, \texttt{\textbackslash ms}, \texttt{\textbackslash no}, \texttt{\textbackslash pf}, \texttt{\textbackslash s}), with extra wrapper packages (bm, bbm, dsfont, xspace, cancel, mathtools). The delivered numbering is the unit's own. Assets: no retained span contains a figure environment, a graphics-inclusion command, a PDF-page inclusion, a table or a TikZ picture; no image file is copied into this package, the package declares no asset dependency and no image file is redistributed with it. Licence: version arXiv:2110.13137v7 of this article is distributed under the Creative Commons Attribution 4.0 International licence, as recorded in the arXiv licence metadata for this exact version and shown on the abstract page https://arxiv.org/abs/2110.13137v7. A full rights notice is delivered at licenses/arxiv-2110.13137-CC-BY-4.0.txt. Characters: every retained excerpt is pure ASCII, so the delivered engine, XeLaTeX with its default Latin Modern text font, reports no missing character.
	\begin{align}\label{eqcms}
		\cms_g\phi=\max_{q\in M}\max_{\substack{P\s T_qM\\ \dim P=d}}\phi(P).
	\end{align}

\begin{lem}\label{lemprod}
Suppose the $d$-dimensional integral current $T$ is area-minimizing (in integral or mod $v$ homology) on a compact  Riemannian manifold $V$. Let $W$ be another  compact Riemannian manifold with $p$ a point in $W$. Define a map $$i_p:V\to V\times W,$$ by
\begin{align*}
i_p(q)=(q,p).
\end{align*}
Then the pushforward current 
\begin{align*}
(i_p)\pf(T)
\end{align*}
is area-minimizing (in integral or mod $v$ homology) on $V\times W$. Furthermore, if $T$ is calibrated by a smooth form $$\phi$$ on $V,$ then $(i_p)\pf T$ is calibrated by a smooth form
\begin{align*}
\pi_V\du\phi,
\end{align*}on $V\times W,$ where $\pi_V$ is the canonical projection of $V\times W$ onto the $V$ factor.
\end{lem}

\begin{proof}
Let us first note that both $\pi_V$ and $i_q\circ\pi_V$ are area-non-increasing. In other words, if $\nu$ is a simple $d$-vector on the tangent space to $V\times W,$ then we have 
\begin{align}
&\no{(\pi_V)\pf \nu}\le \no{\nu},\label{pman}.% \\&\no{(i_q\circ\pi_V)\pf \nu}\le \no{\nu}\label{pqan}. 
\end{align}
Inequality (\ref{pman}) follows directly from the definition of Riemannian length of $d$-vectors and Cauchy-Binet formula. %Inequality (\ref{pqan}) follows from (\ref{pman}) and the fact that $i_p$ is a Riemannian isometric embedding.

Now let $Z$ be an integral current homologous to $(i_p)\pf T$ in $V\times W.$ Then $(\pi_V)\pf Z$ is homologous to $( \pi_V\circ i_p)\pf T$ on $V$.  By definition, $ \pi_V\circ i_p$ is the identity map on $V.$ Thus, $(\pi_V)\pf Z$ is homologous to $T.$ By the area-minimality of $T,$ we have
\begin{align*}
\ms(T)\le \ms((\pi_V)\pf Z).
\end{align*}
By (\ref{pman}), we have
\begin{align*}
\ms((\pi_V)\pf Z)\le\ms(Z).
\end{align*}
Finally, since $i_p$ is a Riemannian isometry, we deduce that
\begin{align*}
\ms((i_p)\pf T)=\ms(T)\le \ms((\pi_V)\pf Z)\le\ms(Z).
\end{align*}
Thus, $(i_p)\pf T$ is area-minimizing.

Direct calculation using (\ref{eqcms}) and (\ref{pman}) shows that $\pi_V\du \phi$ has comass at most that of $\phi$ and $(\pi_V)\du\phi$ calibrates $(i_p)\pf T.$ 
\end{proof}

\end{document}
